%%%%%%%%%%%%%%%%%%%%%%%%%%%%%%%%%%%%%%%%%%%%%%%%%%%%%%%
% % Lines starting with % are comments, which are ignored.
% % This is a handy way of indicating the date and version of
% % your document, to wit:
% %
% % LaTeX sample file
% % Modified March, 2002
% %
%%%%%%%%%%%%%%%%%%%%%%%%%%%%%%%%%%%%%%%%%%%%%%%%%%%%%%%
\DocumentMetadata{
  % PDF standards: ONE key whose value is a comma-separated list --
  % UA-2 (accessibility) plus A-4f (archival; the archival companion of
  % UA-2). Do NOT use separate pdfstandard=... keys: with ordinary
  % key-value parsing only the last one takes effect.
  pdfstandard={UA-2,A-4f},
  lang=en-US,
  tagging=on,
  tagging-setup={
    math/setup=mathml-SE
  }
}

\documentclass{ltx-talk}
% lecture-macros.tex -- shared package setup for lecture-template.tex
% (ltx-talk slides / handouts).
%
% NOTE ON PATHS: every template in this repository is stored one level up
% from where it is compiled (see README.md). All relative \input paths
% below are resolved against the folder the document is compiled in,
% which sits one directory above this file.

% --- Tables, bookmarks, strikethrough ------------------------------------
% booktabs: professional table rules (\toprule, \midrule, \bottomrule)
% bookmark: PDF outline entries (PDF bookmarks)
% ulem[normalem]: \sout strikethrough that survives math mode
\usepackage{booktabs}
\usepackage{bookmark}
\usepackage[normalem]{ulem}
\usepackage{multicol}

% --- Math ------------------------------------------------------------
% amsmath: equation environments and numbering
% mathrsfs: \mathfrak  |  stmaryrd: extra symbols
\usepackage{amsmath}
\usepackage{mathrsfs}
\usepackage{stmaryrd}

% --- Layout -------------------------------------------------------------
% float: [H] float placement
% fontspec: Unicode/OpenType fonts (compile with XeLaTeX or LuaLaTeX)
\usepackage{float}
\usepackage{fontspec}

% --- Colors ---------------------------------------------------------------
% dvipsnames: extended named color palette (the Accessible* colors below
% are built on it);
% table: \rowcolors / \cellcolor for shading tables.
% NOTE: xcolor may only be loaded ONCE per document -- a second
% \usepackage{xcolor} with different options is a fatal "Option clash"
% error, so ALL of its options are listed here.
\usepackage[dvipsnames,table]{xcolor}

% --- Multimedia -----------------------------------------------------------
% animate: embeds audio/video (and legacy Flash) media in a frame.
\usepackage{animate}

% --- Bibliography ----------------------------------------------------------
% biblatex (biber backend): citations for lecture notes.
\usepackage[
  backend=biber,
  style=alphabetic,
  url=true,
  doi=true,
  isbn=true,
  backref=true
]{biblatex}

% --- Colored boxes ----------------------------------------------------------
% talkthemetcolorbox: the colored-box theme for ltx-talk slides
% (shadow/rounded/shaded blocks, showtitle, titlepage styles).
\usepackage[
  shadow,
  rounded,
  shaded,
  blocks,
  showtitle,
  titlepage
]{talkthemetcolorbox}

% --- "theindex" helper environment ------------------------------------------
% Provides a lightweight theindex environment (an itemize) and the
% \@idxitem/\subitem/\subsubitem/\indexspace macros that LaTeX's standard
% theindex machinery expects, so hand-written index pages indent correctly.
\makeatletter
\newenvironment{theindex}
  {\begin{itemize}}
  {\end{itemize}}
\providecommand\@idxitem{\par\hangindent 40\p@}
\providecommand\subitem{\@idxitem \hspace*{20\p@}}
\providecommand\subsubitem{\@idxitem \hspace*{30\p@}}
\providecommand\indexspace{\par \vskip 10\p@ \@plus5\p@ \@minus3\p@\relax}
\makeatother

% --- Index -----------------------------------------------------------------
% imakeidx: builds the printed index.
% The redefinition below appends "|hyperpage" to every index entry so that,
% in the PDF, each index entry is a link to its page.
\usepackage{imakeidx}
\makeatletter
\let\imki@orig@wrindexentrysplit\imki@wrindexentrysplit
\renewcommand\imki@wrindexentrysplit[3]{%
  \imki@orig@wrindexentrysplit{#1}{#2|hyperpage}{#3}%
}

\makeatletter
% --- Section titles: capture + one bookmark per logical section -------------
% \section is redefined so that:
%   1. the section title is captured into the plain-TeX register
%      \mytalksectitletoks (reprinted verbatim on the auto divider frame,
%      see the cmd/section/after hook below -- plain-TeX storage is
%      catcode-safe, unlike expanding the title in a hook), and
%   2. exactly ONE depth-1 PDF bookmark is created per \section call
%      (hyperref's own bookmark is suppressed via \Hy@writebookmark).
\DeclareCommandCopy{\OriginalSectionForBookmarks}{\section}
\newtoks\mytalksectitletoks

\RenewDocumentCommand{\section}{s o m}{%
  \begingroup
    \def\Hy@writebookmark##1##2##3##4##5{}%
    \global\mytalksectitletoks={#3}%
    \IfBooleanTF{#1}
      {\phantomsection\IfNoValueTF{#2}{\OriginalSectionForBookmarks*{#3}}{\OriginalSectionForBookmarks*[#2]{#3}}}%
      {\IfNoValueTF{#2}{\OriginalSectionForBookmarks{#3}}{\OriginalSectionForBookmarks[#2]{#3}}}%
  \endgroup
  \IfBooleanTF{#1}
    {\IfNoValueTF{#2}{\pdfbookmark[1]{#3}{section-star.\thepage.\number\inputlineno}}{\pdfbookmark[1]{#2}{section-star.\thepage.\number\inputlineno}}}%
    {\IfNoValueTF{#2}{\pdfbookmark[1]{Section \thesection: #3}{section.\thesection}}{\pdfbookmark[1]{Section \thesection: #2}{section.\thesection}}}%
}
\makeatother

% --- One PDF bookmark per logical frame (NOT per overlay) ---------------------
% ltx-talk writes one PDF page per overlay, so a naive frametitle hook would
% register a new bookmark for every overlay of the same frame. The
% \g_slideutils_frame_bookmark_int register remembers the last frame that
% got a bookmark, so subsequent overlays of the same frame are skipped.
% This is also what keeps BeamerPresenter's cumulative annotations working:
% one logical frame must stay one logical unit in the PDF outline.
\ExplSyntaxOn
\int_new:N \g_slideutils_frame_bookmark_int

\AddToHook{cmd/frametitle/after}{%
  \tl_if_empty:NF \g__talk_frame_title_tl
    {
      \ifnum\value{frame}=\g_slideutils_frame_bookmark_int\relax
      \else
        \int_gset:Nn \g_slideutils_frame_bookmark_int { \value{frame} }%
        \pdfbookmark[2]{\tl_use:N \g__talk_frame_title_tl}
          {frame.\int_use:N \g_slideutils_frame_bookmark_int}
      \fi
    }
}
\ExplSyntaxOff
\makeatother

% --- Printed index + theme for the colored boxes -------------------------------
% \makeindex[...]: two-column index titled "Index", listed in the TOC.
% The two EditInstance calls style the standard header/footer instances
% of talkthemetcolorbox (cyan background; footer order title/author/
% date/framenumber).
\makeindex[columns=2, title=Index, intoc]\EditInstance{header}{std}{
  background-color = cyan,
  color            = black,
  height           = 0.75cm
}
\EditInstance{footer}{std}{
  background-color = cyan,
  element-order = {title, author, date, framenumber},
}

% --- Auto section divider frame ---------------------------------------------------
% After each \section, insert a centered, large, structure-colored frame
% that prints the section title (from \mytalksectitletoks).
\AddToHook{cmd/section/after}{%
  \begin{frame}
    \centering
    \color{structure}
    \LARGE
    \the\mytalksectitletoks
  \end{frame}%
}

% --- Tagging sockets (tagged PDF structure) ---------------------------------------
% ltx-talk exposes tagging "sockets" where the PDF structure can attach
% section titles; register one and leave it unassigned (none) so the
% class's own tagging keeps control.
\ExplSyntaxOn
\AtBeginDocument{
  \NewTaggingSocketPlug { talk / sec / title } { none } { }
  \AssignTaggingSocketPlug { talk / sec / title } { none }
}
\ExplSyntaxOff

% --- Accessibility roles ----------------------------------------------------------------
% Map \frametitle to the H1 heading role in the tagged PDF, so screen
% readers treat slide titles as headings.
\tagpdfsetup{role / new-tag = frametitle / H1}

\makeatletter
% --- Page label = frame number -------------------------------------------------------
% Redefine what feeds the PDF page label to track the frame counter,
% not the physical page counter. Replace \theframenumber below with
% whatever internal macro/property ltx-talk actually uses to produce
% the "framenumber" footer element — confirm via ltx-talk-frame.dtx.
\renewcommand\thepage{\arabic{frame}}
\makeatother

\addbibresource{refs.bib}

% --- "Problem Time!" slides ------------------------------------------------------------------
% \ProblemSlide{problem}{solution} emits two frames: the problem, then
% its solution after a \pause, each numbered by the problem counter.
\newcounter{problemcounter}
\setcounter{problemcounter}{0}
\newcommand{\ProblemSlide}[2]{
  \stepcounter{problemcounter}
  \begin{frame}
    \frametitle{Problem Time!}
    Problem \theproblemcounter: #1
  \end{frame}

  \begin{frame}
    \frametitle{Problem Time!}
    Solution to Problem \theproblemcounter: \pause

    #2
  \end{frame}
}
% \highlight[<colour>]{<stuff>}
% from https://tex.stackexchange.com/questions/33401/a-version-of-colorbox-that-works-inside-math-environments
\newcommand{\highlight}[2][yellow]{\mathchoice%
  {\colorbox{#1}{$\displaystyle#2$}}%
  {\colorbox{#1}{$\textstyle#2$}}%
  {\colorbox{#1}{$\scriptstyle#2$}}%
  {\colorbox{#1}{$\scriptscriptstyle#2$}}}%

% --- Fonts ------------------------------------------------------------------------------
% Lato (text) + LeteSansMath (math): open, accessible typefaces.
\setmainfont{Lato}[
  Ligatures=TeX,
  Numbers={Lining,Monospaced}
]
\setsansfont{Lato}[
  Ligatures=TeX,
  Numbers={Lining,Monospaced}
]
\setmathfont{LeteSansMath.otf}\usepackage{emoji}

% --- Accessible colors ---------------------------------------------------------------------
% WCAG-contrast-safe variants of red/green for text on white:
% red!80!black and green!40!black both clear the 4.5:1 requirement.
% AltYellow: soft background for highlight boxes.
\colorlet{AccessibleRed}{red!80!black}
\colorlet{AccessibleGreen}{green!40!black}
\definecolor{AltYellow}{RGB}{255,245,170}

% --- Links ---------------------------------------------------------------------------------
% hyperref + link colors: every link uses AccessibleGreen (7.2:1 on
% white), so internal, citation, and URL links all pass WCAG 1.4.3.
\usepackage{hyperref}
\hypersetup{
  colorlinks=true,
  linkcolor=AccessibleGreen,
  urlcolor=AccessibleGreen,
  citecolor=AccessibleGreen,
}

% --- Shared course macros ---------------------------------------------------------------------
% Root macros.tex: math shorthands used across all course documents
% (\RR, \NN, \ZZ, \QQ, \grad, \curl, \tr, \der, \pder, ...).
\newcommand{\RR}{\mathbb{R}}
\newcommand{\NN}{\mathbb{N}}
\newcommand{\ZZ}{\mathbb{Z}}
\newcommand{\QQ}{\mathbb{Q}}
\newcommand{\vv}{\vec{v}}
\newcommand{\ww}{\vec{w}}
\newcommand{\xx}{\vec{x}}
\newcommand{\yy}{\vec{y}}
\newcommand{\zz}{\vec{z}}
\newcommand{\uu}{\vec{u}}
\newcommand{\basis}[1]{\vec{e_{#1}}}
\newcommand{\id}{\mathrm{id}}
\newcommand{\CC}{\mathbb{C}}
\renewcommand{\div}{\operatorname{div}}
\newcommand{\grad}{\operatorname{grad}}
\newcommand{\curl}{\operatorname{curl}}
\newcommand{\tr}{\operatorname{tr}}
\newcommand{\der}[2]{\dfrac{d #1}{d #2}}
\newcommand{\pder}[3][]{\dfrac{\partial^{#1} #2}{\partial #3}}


% % Title and author(s)
%%%%%%%%%%%%%%%%%%%%%%%%%%%%%%%%%%%%%%%%%%%%%%%%%%%%%%%
\title{Lecture 9 - Linearization and Derivative Rules}
\author{Garrett Credi}
\date{10/08/2026}
%%%%%%%%%%%%%%%%%%%%%%%%%%%%%%%%%%%%%%%%%%%%%%%%%%%%%%%
% %
% % The next command allows your in import encapsulated
% % postscript files, .epsf or .eps files, which
% % contain vector graphic image data.
% %
%%%%%%%%%%%%%%%%%%%%%%%%%%%%%%%%%%%%%%%%%%%%%%%%%%%%%%%
% %   The first thanks indicates your affiliation
% %
% %  Just the name here.
% %
% % Your mailing address goes at the end.
% %
% % \thanks is also how you indicate grant support
% %
%%%%%%%%%%%%%%%%%%%%%%%%%%%%%%%%%%%%%%%%%%%%%%%%%%%%%%%

%%%%%%%%%%%%%%%%%%%%%%%%%%%%%%%%%%%%%%%%%%%%%%%%%%%%%%%
% $ The next command ensures the margins are 1 inch on
% % all sides.
% %
%%%%%%%%%%%%%%%%%%%%%%%%%%%%%%%%%%%%%%%%%%%%%%%%%%%%%%%
\begin{document}
\begin{frame}
  \maketitle
\end{frame}

\begin{frame}
  \frametitle{Post Exam Recap} 

  \begin{enumerate}
    \item Solutions to Exam 1 are up on Canvas, I'm happy to answer any questions about it if people want
    \item Exams should be graded by tonight or tomorrow night, with curve details to follow after I release the scores
    \item One important note so far: follow the instructions in the questions! Many people did not include reasoning for parts b and c of question 2; this will decrease your scores!
  \end{enumerate}
\end{frame}

\begin{frame}
  \frametitle{Quiz Time!}

  Let's get into neat rows and spread out for the quiz! \pause

  Just kidding :) no quiz today
\end{frame}

\section{Linearization}

\begin{frame}
  \frametitle{Best Linear Approximation}

  Recall that the \textit{derivative} of a function $f(x)$
  at a point $x = a$
  was equal to the slope of the \textit{best linear approximation}
  to $f(x)$ at $x = a$. \pause

  \begin{equation}
    x \approx a \quad \implies \quad f(x) \approx f(a) + f'(a)(x - a).
  \end{equation} \pause

  \begin{block}{Linearization}\index{linearization}\index{derivative!linearization}
    The linearization of a differentiable function $f(x)$ at a point $x = a$
    is the best linear approximation to $f(x)$ at $x = a$ and is equal to \pause
    \begin{equation}
      \mathscr{L}_{x = a}(x) = f(a) + f'(a)(x - a)
    \end{equation}
  \end{block}
\end{frame}

\begin{frame}
  \frametitle{Using the Linearization to Approximate I}

  Using linearizations gives us \textit{the strongest} ways
  approximate functions in a quick and accurate way. \pause

  \textbf{Example 1:} The \textit{logistic population growth} model
  is a mathematical description of how the size of a population $P(t)$ grows
  over some time period $t$ with some growth factor $k$ and some carrying capacity $N$. It assumes that:
  \begin{equation}
    \dfrac{dP}{dt} = k P(N - P)
  \end{equation}
  Suppose that we have a population $P(t)$, where $t$ is measured in days, which: \pause
  \begin{enumerate}
    \item Starts with $1,000$ individuals \pause So $P(0) = 1000$ \pause
    \item Has a growth rate of $0.0001$ \pause So $k = 0.0001$ \pause
    \item Has a carrying capacity of $11,000$ individuals \pause So $N = 11000$ \pause
  \end{enumerate}

  What is our best estimate for the population size one year after we started measuring?
\end{frame}

\begin{frame}
  \frametitle{Using the Linearization to Approximate II}
  \begin{equation}
    \dfrac{dP}{dt} = kP(N - P), \quad P(0) = 1000, \quad k = 0.0001, \quad N = 11000, \pause \quad P(1) \approx ?
  \end{equation}

  Let's linearly approximate! The linearization of $P(t)$ around $t = 0$ is: \pause
  \begin{equation}
    \mathscr{L}_{t = 0}(t) = P(0) + P'(0)(t - 0)
  \end{equation} \pause

  We know $P(0) = 1000$ \pause what about $P'(0)$? \pause We use the $\dfrac{dP}{dt}$ equation! \pause

  \begin{align}
    P'(0) &= \dfrac{dP}{dt}\mid_{t = 0} = k \cdot P(0) \cdot (N - P(0)) \\ \pause
          &=  0.0001 \cdot 1000 \cdot (11000 - 1000) = \dfrac{1}{10} \cdot 10000 \\ \pause
          &= 1000 
  \end{align} \pause
  And so: \pause

  \begin{equation}
    P(1) \approx P(0) + P'(0)(1 - 0) = 1000 + 1000 \cdot 1 = 2000
  \end{equation}
\end{frame}

\begin{frame}
  \frametitle{Using the Linearization to Approximate III}

  Is this an over-approximation ($\mathscr{L}_{t = 0}(1) > P(1)$) or an under-approximation ($\mathscr{L}_{t = 0}(1) < P(1)$)? What would we need to determine that? \pause

  Suppose we know that $P''(0) = 900$. Does that mean our linearization gave an over-approximation or an under-approximation? \pause

  Since $P''(0) > 0$, $P(t)$ is \pause \textit{concave-up} at $t = 0$,
  which means our approximation is an \textit{under} approximation! \pause

  (Concave up => graph is above its tangent line)
\end{frame}

\begin{frame}
  \frametitle{Using the Linearization to Approximate IV}

  We can also approximate functions that are:
  \begin{multicols}{2}
    \begin{enumerate}
      \item Easy to define but \pause
      \item Hard to evaluate
    \end{enumerate}
  \end{multicols}
  like $f(x) = \sqrt{x}$, around known values. \pause

  \textbf{Example 2:} Approximate $\sqrt{25.1}$. \pause

  Approximation should always have us think of linearization: what do we need to linearize? \pause

  We need:
  \begin{enumerate}
    \item An initial value $x = a$ around which to approximate where \pause
    \item We know $f(a)$ AND \pause
    \item We know $f'(a)$. \pause
  \end{enumerate}

  What is a value $a \approx 25.1$ where we know $\sqrt{a}$ and the derivative of the sqrt function at $x = a$?
\end{frame}

\begin{frame}
  \frametitle{Using the Linearization to Approximate V}

  We know that: \pause
  \begin{equation}
    \dfrac{d}{dx} \sqrt{x} = \dfrac{1}{2\sqrt{x}}
  \end{equation} \pause

  so knowing $\dfrac{d}{dx}\mid_{x = a} \sqrt{x}$ is as hard as knowing $\sqrt{a}$. What should $a$ be? \pause

  We should set $a = 25$ since $25 \approx 25.1$ AND $\sqrt{25} = 5$ is easy!

  Then we get:
  \begin{align}
    \sqrt{25.1} &\approx \sqrt{25} + \dfrac{1}{2 \cdot \sqrt{25}} (25.1 - 25) \\
                &= 5 + \dfrac{1}{2 \cdot 5} \cdot (0.1) \\
                &= 5 + \dfrac{1}{10} \cdot 0.1 \\
                &= 5 + 0.01 = 5.01
  \end{align}
\end{frame}

\begin{frame}
  \frametitle{Using the Linearization to Approximate VI}
  Let's say we also knew that:
  \begin{equation}
    \dfrac{d}{dx} \dfrac{1}{2 \sqrt{x}} = \dfrac{-1}{4 \sqrt{x^3}}. 
  \end{equation} \pause

  Is our approximation then an over or under approximation? \pause

  Well:
  \begin{align}
    \dfrac{d^2}{dx^2}\mid_{x = 5} \sqrt{x} &= \left(-\dfrac{1}{4\sqrt{x^3}}\right)\mid_{x = 5} \\
                                           &= -\dfrac{1}{4\sqrt{125}} < 0
  \end{align}

  So this is an \textit{over} approximation since we are concave down at $x = 5$.
\end{frame}

\section{Derivative Rules}

\begin{frame}
  \frametitle{Limit Definition is Hard}

  How can we get that $\dfrac{d}{dx} \dfrac{1}{2 \sqrt{x}} = \dfrac{-1}{4 \sqrt{x^3}}$? \pause

  \begin{align}
    \dfrac{d}{dx} \dfrac{1}{2\sqrt{x}} &= \lim\limits_{h \rightarrow 0} \dfrac{\dfrac{1}{2\sqrt{x + h}} - \dfrac{1}{2\sqrt{x}}}{h} \\ \pause
                                       &= \dfrac{1}{2}\lim\limits_{h \rightarrow 0} \dfrac{\dfrac{\sqrt{x + h}}{x + h} - \dfrac{\sqrt{x}}{x}}{h} \\ \pause
                                       &= \dfrac{1}{2} \lim\limits_{h \rightarrow 0} \dfrac{\dfrac{x\sqrt{x + h}}{x(x + h)} - \dfrac{(x + h)\sqrt{x}}{x(x + h)}}{h} \\ \pause
                                       &= \dfrac{1}{2} \lim\limits_{h \rightarrow 0} \dfrac{\sqrt{x^3 + hx^2} - \sqrt{(x + h)^2 \cdot x}}{h x (x + h)} \\ \pause
                                       &=  \ldots \pause
  \end{align}

  GROSS!!
\end{frame}

\begin{frame}
  \frametitle{Derivative Rules Rule!}

  We'll see that a much easier way to compute this derivative
  involved using \textit{derivative rules} - variants of the sum/product/quotient limit rules we discussed a couple weeks ago. \pause

  For a sneak peek: \pause
  \begin{align}
    \dfrac{d}{dx} \dfrac{1}{2\sqrt{x}} &= \dfrac{d}{dx} \dfrac{1}{2} \dfrac{1}{x^{\frac{1}{2}}} \\
                                       &= \dfrac{d}{dx} \dfrac{1}{2} x^{-\frac{1}{2}} \\
                                       &= \dfrac{1}{2} \dfrac{d}{dx} x^{-\frac{1}{2}} \\
                                       &= \dfrac{1}{2} \cdot -\dfrac{1}{2} \cdot x^{-\frac{1}{2} - 1} \\
                                       &= \dfrac{-1}{4\sqrt{x^3}}
  \end{align}
\end{frame}

\begin{frame}
  \frametitle{Derivative Rules}

  These will provide a \textit{synthetic} approach to computing derivatives, and are split up similar to our rules for limits: \pause some rules tell us how our operation acts on ``simple" things, \pause the rest tell us how our operation acts on ``combinations" of things.

  Atomic Rules: \pause
  \begin{multicols}{2}
    \begin{itemize}
      \item Power Rule: $\dfrac{d}{dx} x^n$ 
      \item Trigonometric Rules: $\dfrac{d}{dx} \sin(x)$ etc.
      \item Exponential/Logarithm Rules: $\dfrac{d}{dx} e^x$, $\dfrac{d}{dx} \ln(x)$
    \end{itemize}
  \end{multicols}

  Combination Rules: \pause
  \begin{multicols}{2}
    \begin{itemize}
      \item Sum/Difference Rule: $\dfrac{d}{dx} (f(x) \pm g(x))$ 
      \item Product Rule: $\dfrac{d}{dx} f(x)g(x)$ 
      \item Chain Rule: $\dfrac{d}{dx} g(f(x))$ 
      \item Constant Multiple Rule: $\dfrac{d}{dx} c \cdot f(x)$
      \item Quotient Rule: $\dfrac{d}{dx} \dfrac{f(x)}{g(x)}$ 
      \item Inverse Function Rule: $\dfrac{d}{dx} f^{-1}(x)$
    \end{itemize}
  \end{multicols}
\end{frame}

\begin{frame}
  \frametitle{Sum/Difference Rule}

  Perhaps the most straightforward derivative rule:
  \begin{block}{Sum/Difference Rule}\index{derivative!sum/difference rule}
    The derivative of a sum/difference of functions
    is the sum/difference (resp.) of the derivatives of the functions: \pause
    \begin{equation}
      \dfrac{d}{dx} (f(x) + g(x)) = \left(\dfrac{d}{dx} f(x)\right) + \left(\dfrac{d}{dx} g(x)\right).
    \end{equation} \pause
  \end{block}

  Why? \pause What is the slope of a sum of two linear functions? \pause
  \begin{equation}
    (f(a) + f'(a)(x - a)) + (g(a) g'(a)(x - a)) = (f(a) + g(a)) + (f'(a) + g'(a))(x - a)
  \end{equation}
\end{frame}

\begin{frame}
  \frametitle{Constant Multiple Rule}

  Similarly: 
  \begin{block}{Constant Multiple Rule}
    The derivative of a function multiplied by a constant
    is that constant multiplied by the derivative of the function: \pause
    \begin{equation}
      \dfrac{d}{dx}(c \cdot f(x)) = c \dfrac{d}{dx} f(x)
    \end{equation} \pause
  \end{block}

  Why? \pause What is the slope of a constant multiple of a linear function? \pause
  \begin{equation}
    c \cdot (f(a) + f'(a) (x - a)) =  (c \cdot f(a)) + (c \cdot f'(a))(x - a)
  \end{equation} \pause

  \begin{alertblock}{WARNING!}
    The factor must actually be a constant!
    \begin{align}
      \dfrac{d}{dx} x^2 &= 2x 
                        \neq x \cdot \dfrac{d}{dx} x = x \cdot 1 = x
    \end{align}
  \end{alertblock}
\end{frame}

\begin{frame}
  \frametitle{Example}
  We can use this to compute derivatives of functions like:
  \begin{equation}
    f(x) = 5x^3 + 3x^2 + 2x + 5
  \end{equation} \pause

  Since we already know that
  \begin{equation}
    \dfrac{d}{dx} x^3 = 3x^2, \quad \dfrac{d}{dx} x^2 = 2x, \quad \dfrac{d}{dx} x = 1, \quad \dfrac{d}{dx} 1 = 0
  \end{equation} \pause

  we get: \pause
  \begin{align}
    f'(x) &= \dfrac{d}{dx}(5x^3 + 3x^2 + 2x + 5) \\ \pause
          &= \dfrac{d}{dx} 5x^3 + \dfrac{d}{dx} 3x^2 + \dfrac{d}{dx} 2x + \dfrac{d}{dx} 5 \\ \pause
          &= \dfrac{d}{dx} (5 \cdot x^3) + \dfrac{d}{dx}(3 \cdot x^2) + \dfrac{d}{dx} (2 \cdot x) + \dfrac{d}{dx}(5 \cdot 1) \\ \pause
          &= 5 \cdot \dfrac{d}{dx} x^3 + 3 \cdot \dfrac{d}{dx} x^2 + 2 \cdot \dfrac{d}{dx} x + 5 \cdot \dfrac{d}{dx} 1 \\ \pause
          &= 5 \cdot (3x^2) + 3 \cdot (2x) + 2 \cdot 1 + 5 \cdot 0 \\ \pause
          &= 15x^2 + 6x + 2
  \end{align}
\end{frame}

\begin{frame}
  \frametitle{Power Rule I}

  What about for more complicated polynomials \pause e.g. what if the degree is $> 3$? \pause

  We know that:
  \begin{align}
    \dfrac{d}{dx} x^3 &= 3x^2 \\
    \dfrac{d}{dx} x^2 &= 2x \\
    \dfrac{d}{dx} x^1 &= 1 \\
    \dfrac{d}{dx} x^0 &= 0
  \end{align} \pause

  Can we see a pattern? \pause

  Guess: $\dfrac{d}{dx} x^n = n \cdot x^{n - 1}$.
\end{frame}

\begin{frame}
  \frametitle{Power Rule II}

  Why?

  \begin{align}
    \dfrac{d}{dx} x^n &= \lim\limits_{h \rightarrow 0} \dfrac{(x + h)^n - x^n}{h} \\ \pause
                      &= \lim\limits_{h \rightarrow 0} \dfrac{\_ x^n + \_ x^{n-1} \cdot h + \_ x^{n-2} h^2 + \ldots + \_ x^0 h^n - x^n}{h} \\ \pause
                      &= \lim\limits_{h \rightarrow 0} (\_ x^{n - 1} \cdot 1 + \_ x^{n - 1} h + \ldots + \_ x^0 h^{n - 1}) \\ \pause
                      &= \_ x^{n-1}
  \end{align} \pause

  What is the coefficient on $hx^{n-1}$ in $(x + h)^n$? \pause It's $n$! \pause

  So we know that $\dfrac{d}{dx} x^n = n x^{n - 1}$!
\end{frame}

\begin{frame}
  \frametitle{Power Rule III}

  \begin{block}{Power Rule}\index{derivative!power rule}
    \begin{equation}
      \dfrac{d}{dx} x^n = nx^{n-1}
    \end{equation}
  \end{block} \pause

  E.g.: \pause
  \begin{itemize}
    \item $\dfrac{d}{dx} x^{100} = $ \pause
    \item $\dfrac{d}{dx} x^{67} = $ \pause
    \item $\dfrac{d}{dx} x$ 
  \end{itemize}
\end{frame}

\begin{frame}
  \frametitle{Power Rule IV}

  Note: We technically only showed the limit definition worked for \textit{integer} powers ($x^1, x^2, x^3$) \pause as opposed to fractional powers, i.e. roots ($x^{\frac{3}{2}}, \sqrt{x}, \sqrt[3]{x^2}$). \pause
  However, with product rule or inverse rule it's easy enough to show that power rule also works for fractional, and even all real-valued power functions. \pause

  E.g. $\dfrac{d}{dx}x^{\pi} = $ \pause

  \begin{alertblock}{WARNING!!}
    Power rule only applies to \textit{\textbf{POWER FUNCTIONS}} \pause
    i.e. $f(x) = x^c$ for some constant $c$.

    It does NOT apply to \textit{exponential functions} like $f(x) = c^x$!
    They may look similar, but their derivatives are VASTLY different!
  \end{alertblock}
\end{frame}

\begin{frame}
  \frametitle{Exponential Rule I}

  What is the derivative of an exponential function? \pause

  \begin{align}
    \dfrac{d}{dx} c^x &= \lim_{h \rightarrow 0} \dfrac{c^{x + h} - c^x}{h}  \\ \pause
                      &= \lim_{h \rightarrow 0} c^x \dfrac{c^h - 1}{h} \\ \pause
                      &= \left(\lim\limits_{h \rightarrow 0} \dfrac{c^h - 1}{h}\right) c^x
  \end{align} \pause

  it is just some constant multiple of the original function! \pause

  What is the constant factor $\lim\limits_{h \rightarrow 0} \dfrac{c^h - 1}{h}$?
\end{frame}

\begin{frame}
  \frametitle{Exponential Rule II}

  Define a constant $e$ so that $\lim\limits_{h \rightarrow 0} \dfrac{e^h - 1}{h} = 1$ (i.e. $\dfrac{d}{dx} e^x = e^x$). Then:

  \begin{align}
    \lim\limits_{h \rightarrow 0} \dfrac{c^h - 1}{h} &= \pause \lim\limits_{h \rightarrow 0} \dfrac{(e^{\ln(c)})^h - 1}{h} \\ \pause
                                                     &= \lim\limits_{h \rightarrow 0} \dfrac{e^{\ln(c) h} - 1}{h} \\ \pause
                                                     &= \lim\limits_{h \rightarrow 0} \dfrac{\ln(c)}{\ln(c)} \dfrac{e^{\ln(c) h} - 1}{h} \\ \pause
                                                     &= \ln(c) \lim\limits_{h \rightarrow 0} \dfrac{e^{\ln(c) h} - 1}{\ln(c) h} \\ \pause
                                                     &= \ln(c) \lim\limits_{z \rightarrow 0} \dfrac{e^z - 1}{z} = \ln(c) \cdot 1 = \ln(c)
  \end{align}
\end{frame}

\begin{frame}
  \frametitle{Exponential Rule III}
  This gives us:

  \begin{block}{Exponential Rule}\index{derivative!exponential rule}
    The derivative of an exponential function, i.e. $c^x$ for a constant $c$ is:
    \begin{equation}
      \dfrac{d}{dx} c^x = \ln(c) c^x
    \end{equation}
  \end{block} \pause

  Exponential functions growing rapidly => it's derivative grows very large very quickly.

  Very different behavior to that for power rule $\dfrac{d}{dx} x^c = cx^{c - 1}$.
\end{frame}

\begin{frame}
  \frametitle{Exponential Rule IV}

  What is the $n$'th derivative of an exponential function $c^x$? \pause

  First derivative: $\ln(c) c^x$. \pause

  Second derivative: $\dfrac{d}{dx} (\ln(c) c^x) = \ln(c) \dfrac{d}{dx} c^x = \ln(c)^2 c^x$ \pause

  Third derivative: $\ln(c)^3 c^x$. \pause

  \ldots

  $n$'th derivative: $\ln(c)^n c^x$.
\end{frame}

\begin{frame}
  \frametitle{Worksheet Time!}

  Let's get into groups and work on the worksheet!
\end{frame}

\begin{frame}
  \frametitle{End of Class}

  That's all folks!
\end{frame}
\end{document}
